\scanpage{18}
\mspage{4}
On dit qu'une partie-bord \inserted{$B$ de $X$} est \uline{compatible} avec
une famille de parties $X_i$ de $X$, s'il
existe un voisinage cylindrique qui est
compatible \inserted{(au sens cylindrique)} avec chacun des $X_i$.
\margin{utiliser plutôt la terminologie de \uline{transversalité}.}

Revenons à $Y\subset X$ et un voisinage
conique \inserted{\uuline{strict}} $\bar T$ de $Y$. On dit que $\bar T$ est
\inserted{\uline{strictement}} compatible avec une famille de $X_i\subset X$
si \inserted{elle est compatible avec chaque $X_i$
$(X_i\cap\bar T=I.X_i\cap\dot T-)$, et si}
$\dot T$ est un sous-espace bord de
$X^*=X\setminus\mathring T$, compatible avec la famille des
$X_i^*=X^*\cap X_i$. Donc cela implique que
$X_i\cap\dot T$ est un sous-espace bord (non seulement
de $X_i\cap\bar T$, mais aussi) de $X_i^*$.
\margin{\deleted{\unsure{appeler transversalité ?}}}
\margin{On dira que $X_i$ est
\deleted{\uline{\unsure{strictement}}}\uline{transversal}
\deleted{compatible \unsure{avec $\bar T$}}
\deleted{si $X_i\cap Y=0.\unsure{X_i}$}
\unsure{désir : transversalité}
\deleted{stricte} \unsure{transversalité} :
\[
0.\dot X_i=X_i\cap Y
\]
i.e. $X_i=\overline{X_i\setminus X_i\cap Y}$}

\textcircled{4}. \deleted{Soit} Un voisinage conique (absolu)
$\bar T\deleted{\unsure{_\varepsilon}}$ de
$Y$ est dit \uline{équisingulier} si $\dot T\longrightarrow Y$ est une
fibration loc. triviale. Il est dit \uline{équisingulier
pour une famille de} $\bar T_i\subset\bar T$, si \uline{de plus}
les $\unsure{Y_i|T_i}$ sont coniques, et $\dot T_i\longrightarrow Y$ sont des
fibrations localement triviales --- et \unsure{même}
\margin{17}
